\section{Planteamiento del problema}

\subsection{Contexto internacional}
El sistema Braille se estructura en celdas de dos columnas por tres filas. En
visión por computadora, el reconocimiento óptico de Braille (OBR) pasó de la
umbralización y la morfología matemática \parencite{antonacopoulos2004} al
aprendizaje profundo supervisado \parencite{li2020}. Los detectores de una etapa
localizan la celda completa y la clasifican en 63 categorías, lo que confiere
invarianza ante rotaciones y deformaciones del papel
\parencite{ovodov2021,sun2025}. Sin embargo, esos avances se formularon para
documentos editoriales impresos en inglés, ruso o chino, o dependen de servicios
en la nube.

\subsection{Contexto nacional}
En el Perú, la normativa de educación inclusiva reconoce el derecho de las
personas con discapacidad a la Educación Básica Regular. Pese a ello, el Informe
Defensorial N.º 183 reportó que el Programa Presupuestal 106 representaba apenas
el 0.62\,\% del presupuesto educativo y que más del 88\,\% de las personas con
discapacidad en edad escolar no estaba matriculada \parencite{defensoria2019}.
Los Censos Nacionales 2017 registraron 3\,209\,261 personas con discapacidad
(10.3\,\%), y la dificultad para ver fue la más frecuente (48.3\,\%;
\citeauthor{inei2018}, \citeyear{inei2018}).

\subsection{Realidad local}
En Moquegua, y en particular en Ilo y Mariscal Nieto, los estudiantes con ceguera
asisten a instituciones educativas regulares y reciben apoyo itinerante del CEBE
``Corazón de Jesús''. Su herramienta cotidiana es la regleta y el punzón sobre
cartulinas de 120 a 160 g/m\textsuperscript{2}, que producen relieves de
aproximadamente 0.5 mm sin tinta ni contraste cromático. El docente de aula, que
atiende además a entre 25 y 30 estudiantes videntes, no puede leer el cuaderno;
la corrección queda retenida hasta la visita del especialista SAANEE/SAEE, que
suele espaciarse semanas. La retroalimentación formativa, cuyo valor depende de
la inmediatez \parencite{hattie2007}, queda anulada.

\subsection{Justificación de la solución}
Frente a este escenario, el proyecto diseña y evalúa una red neuronal
convolucional ligera, basada en detectores de una etapa y cuantización INT8,
capaz de reconocer celdas Braille Grado 1 manuscritas y ejecutarse sin conexión
en teléfonos de gama media/baja. El aporte será un corpus anotado de Braille
escolar manuscrito en español peruano y evidencia comparativa sobre el
compromiso precisión--eficiencia--robustez.
